\documentclass[a4paper,12pt]{article}
\usepackage[T2A]{fontenc}
\usepackage[utf8]{inputenc}
\usepackage[russian]{babel}
\usepackage{amsmath,amssymb,mathtools}
\usepackage{hyperref}
\usepackage{booktabs}
\usepackage{geometry}
\geometry{top=2cm, bottom=2cm, left=2.5cm, right=2.5cm}

\title{Конспект лекции 5: \\ Основы обработки естественного языка}
\author{Липкин С.М.}
\date{}

\begin{document}
\maketitle

\section{Что такое NLP?}

Обработка естественного языка (Natural Language Processing, NLP) — область на стыке искусственного интеллекта, лингвистики и информатики, изучающая методы взаимодействия компьютера и человеческого языка.

\textbf{История развития:} зародилась в 1950-х годах (тест Тьюринга). Ранние подходы использовали правила и статистику. Современный этап (2013--настоящее время) характеризуется доминированием глубокого обучения: word2vec, seq2seq, Transformer, LLM.

\textbf{Базовый конвейер NLP:}
\[
\text{Текст} \xrightarrow{\text{Предобработка}} \text{Токены} \xrightarrow{\text{Векторизация}} \text{Эмбеддинги} \xrightarrow{\text{Модель}} \text{Ответ}
\]

\textbf{Области применения:} чат-боты, поиск, машинный перевод, анализ тональности, голосовые помощники, аннотация текстов, вопросно-ответные системы.

\section{Токенизация}

Токенизация — разбиение текста на минимальные осмысленные единицы (токены). Это обязательный первый шаг NLP-конвейера.

\subsection{Основные подходы}

\begin{enumerate}
    \item \textbf{Word-level:} разделение по пробелам и знакам пунктуации. Даёт большой словарь ($10^5$--$10^6$ токенов). Проблема: OOV (out-of-vocabulary) — модель не знает редких слов.
    
    \item \textbf{Subword (подсловный):} разбиение на подслова. Баланс размера словаря и покрытия. Решает OOV. \textbf{BPE (Byte-Pair Encoding)} — итеративное слияние частых пар символов, используется в GPT. \textbf{WordPiece} — слияние по максимизации likelihood, используется в BERT.
    
    \item \textbf{Character-level:} словарь состоит из символов алфавита (30--200 токенов). Минимальный размер словаря, но длинные последовательности (каждое слово — несколько символов), сложнее улавливать смысл.
\end{enumerate}

\section{Нормализация текста}

Набор шагов для приведения текста к единообразному виду:

\subsection{Базовые шаги}

\begin{enumerate}
    \item \textbf{Нижний регистр:} «Москва» и «москва» — один токен. Уменьшает словарь.
    \item \textbf{Удаление стоп-слов:} предлоги, союзы, местоимения несут мало смысла и удаляются (актуально не для всех задач — для некоторых, наоборот, важны).
    \item \textbf{Очистка от спецсимволов:} удаление пунктуации, URL, HTML-тегов, цифрового мусора.
\end{enumerate}

\subsection{Стемминг и лемматизация}

\textbf{Стемминг:} механическое отсечение окончаний по правилам. Быстро, но грубо. Без учёта контекста. Пример: «бегущий» $\to$ «бегущ», «красивее» $\to$ «красив». Алгоритмы: Портера, Snowball. Есть версии для русского языка.

\textbf{Лемматизация:} лингвистически корректное приведение к начальной форме (лемме) через словарь и морфологический анализ. Точнее, но медленнее стемминга. Пример: «бежал» $\to$ «бежать», «красивее» $\to$ «красивый». Инструменты: NLTK, spaCy, pymorphy2 (русский).

\section{Типовые задачи NLP}

\begin{enumerate}
    \item \textbf{Классификация текста:} определение категории (спам, тематика, жанр).
    \item \textbf{Анализ тональности:} определение эмоциональной окраски (позитив/негатив/нейтрально).
    \item \textbf{Машинный перевод:} автоматический перевод между языками.
    \item \textbf{NER (Named Entity Recognition):} извлечение именованных сущностей (имена, даты, организации, локации).
    \item \textbf{Генерация текста:} создание связного текста (саммаризация, диалоги).
    \item \textbf{Вопросно-ответные системы (QA):} поиск ответа на вопрос в заданном контексте.
    \item \textbf{Языковое моделирование:} предсказание следующего слова.
    \item \textbf{Парсинг:} синтаксический анализ, dependency parsing.
    \item \textbf{Textual Entailment:} определение логической связи между текстами.
\end{enumerate}

\section{RNN в NLP}

\subsection{Принцип работы}

RNN обрабатывают элементы последовательности по одному, сохраняя скрытое состояние (память), которое передаётся на каждом шаге:
\[
h_t = f(W h_{t-1} + U x_t + b)
\]

Язык — последовательная структура, поэтому RNN хорошо подходят: учитывают порядок слов и могут моделировать контекст.

\subsection{Проблемы RNN}

При длинных последовательностях базовая RNN теряет информацию о ранних элементах из-за многократного умножения градиентов (исчезающий градиент).

\subsection{LSTM (Long Short-Term Memory)}

Hochreiter \& Schmidhuber, 1997. Решает проблему через gating-механизм и отдельную ячейку состояния (cell state):
\[
\begin{aligned}
f_t &= \sigma(W_f [h_{t-1}, x_t] + b_f) \quad \text{(forget gate)} \\
i_t &= \sigma(W_i [h_{t-1}, x_t] + b_i) \quad \text{(input gate)} \\
\tilde{C}_t &= \tanh(W_C [h_{t-1}, x_t] + b_C) \quad \text{(кандидат)} \\
C_t &= f_t \odot C_{t-1} + i_t \odot \tilde{C}_t \\
o_t &= \sigma(W_o [h_{t-1}, x_t] + b_o) \quad \text{(output gate)} \\
h_t &= o_t \odot \tanh(C_t)
\end{aligned}
\]

Cell state — «конвейер» информации с минимальными изменениями. Forget gate решает, что забыть, input gate — что добавить, output — что вывести.

\subsection{GRU (Gated Recurrent Unit)}

Упрощённая версия LSTM с двумя gate:
\[
\begin{aligned}
r_t &= \sigma(W_r [h_{t-1}, x_t] + b_r) \quad \text{(reset gate)} \\
z_t &= \sigma(W_z [h_{t-1}, x_t] + b_z) \quad \text{(update gate)} \\
\tilde{h}_t &= \tanh(W [r_t \odot h_{t-1}, x_t] + b) \\
h_t &= (1-z_t) \odot h_{t-1} + z_t \odot \tilde{h}_t
\end{aligned}
\]

Меньше параметров, быстрее обучение, меньше памяти, сопоставимое качество.

\subsection{LSTM vs GRU}

\begin{itemize}
    \item LSTM: 3 gate, отдельный cell state, больше параметров, лучше на длинных последовательностях.
    \item GRU: 2 gate, единое скрытое состояние, меньше параметров, быстрее, сопоставимое качество на коротких текстах.
\end{itemize}

\section{Seq2Seq и механизм внимания}

\subsection{Encoder-Decoder}

Архитектура для задач, где вход и выход — последовательности разной длины (перевод, саммаризация):
\begin{itemize}
    \item \textbf{Encoder:} RNN/LSTM читает входную последовательность и генерирует последовательность скрытых состояний $h_1, \dots, h_T$.
    \item \textbf{Decoder:} RNN/LSTM генерирует выходную последовательность, используя контекст от энкодера.
\end{itemize}

\textbf{Проблема:} при длинных предложениях информация теряется при «упаковке» в один вектор контекста.

\subsection{Механизм внимания (Bahdanau, 2015)}

На каждом шаге декодер вычисляет веса $\alpha_{tj}$ для всех состояний энкодера:
\[
e_{tj} = a(s_{t-1}, h_j), \quad
\alpha_{tj} = \frac{\exp(e_{tj})}{\sum_{k} \exp(e_{tk})}, \quad
c_t = \sum_{j=1}^{T} \alpha_{tj} h_j
\]

Контекстный вектор $c_t$ — взвешенная сумма состояний энкодера с акцентом на релевантные части.

\textbf{Значение:} внимание стало фундаментом архитектуры Transformer (Vaswani et al., 2017), который лежит в основе GPT, BERT и всех современных LLM.

\section{Применение RNN в NLP}

\subsection{Классификация текста и тональность}

RNN читает текст слово за словом, накапливая контекст в скрытом состоянии. Последнее скрытое состояние (или max-pooling по всем шагам) подаётся на классификатор (Dense + Softmax).

Пример: отзыв «Этот фильм великолепен!» проходит через 4 шага LSTM, финальное состояние кодирует положительную тональность, Softmax выдаёт P(позитив) = 0.92.

\textbf{BiLSTM:} двунаправленный LSTM обрабатывает текст слева направо и справа налево, захватывая полный контекст.

\subsection{NER: BiLSTM-CRF}

Стандартная архитектура для NER:
\begin{itemize}
    \item \textbf{BiLSTM} извлекает контекстные признаки для каждого токена (с учётом контекста слева и справа).
    \item \textbf{CRF-слой} (Conditional Random Field) обеспечивает согласованность меток: например, B-PER не может идти после I-ORG; I-ORG может идти только после B-ORG.
\end{itemize}

\textbf{BIO-тегирование:} B (begin) — начало сущности, I (inside) — продолжение, O (outside) — не сущность. Типы: PER (персона), ORG (организация), LOC (локация), DATE (дата).

Пример: «Иван работает в Газпром с 2010 года» $\to$ B-PER, O, O, B-ORG, O, DATE, O.

\subsection{Машинный перевод}

Seq2Seq с вниманием — стандарт перевода до появления Transformer. Энкодер читает исходное предложение, декодер генерирует перевод. На каждом шаге внимание сфокусировано на релевантных частях исходного предложения.

Пример: при переводе «Я изучаю нейронные сети» на шаге «study» максимальный вес attention направлен на «изучаю».

\subsection{Языковое моделирование и генерация}

Языковая модель предсказывает вероятность последовательности:
\[
P(w_1, \dots, w_T) = \prod_{t=1}^{T} P(w_t | w_{<t})
\]

На этапе генерации модель работает \textbf{авторегрессивно}: на каждом шаге подаёт собственный выход как следующий вход.

\textbf{Стратегии декодирования:}
\begin{itemize}
    \item Greedy (argmax): выбирается самый вероятный токен.
    \item Temperature sampling: $P(y) = \text{softmax}(\text{logits} / \tau)$. $\tau \to 0$ — greedy, $\tau \to \infty$ — равномерное.
    \item Top-K sampling: семплинг из K лучших токенов.
    \item Top-P (nucleus) sampling: семплинг из минимального набора токенов, покрывающих массу вероятности $p$.
\end{itemize}

\section{Сравнение подходов}

\begin{center}
\begin{tabular}{lccc}
\toprule
\textbf{Критерий} & \textbf{RNN} & \textbf{LSTM/GRU} & \textbf{Transformer} \\
\midrule
Длинные зависимости & Слабо & Хорошо & Отлично \\
Параллелизация & Нет & Нет & Да \\
Скорость обучения & Быстрая & Средняя & Высокая (GPU) \\
Качество (SOTA) & Низкое & Среднее & SOTA \\
Потребление памяти & Мало & Средне & Много \\
Типичное применение & Простые задачи & Перевод, NER & LLM, чат-боты \\
\bottomrule
\end{tabular}
\end{center}

\section{Заключение}

NLP прошёл путь от правил и статистики до глубокого обучения. RNN/LSTM/GRU остаются актуальны для задач с ограниченными ресурсами, streaming-данных и real-time обработки. Transformer доминирует в качестве, но требует больше вычислительных ресурсов.

Ключевые этапы конвейера — токенизация и нормализация — критически важны для качества любой NLP-модели. Выбор подхода (word, subword, character) зависит от языка, задачи и доступных ресурсов.

\end{document}
